\documentclass[11pt]{article}
\usepackage{fontspec}
\setmainfont{Times New Roman}
\setsansfont{Arial}
\setmonofont{Consolas}
\usepackage{amsmath,amssymb}
\usepackage{geometry}
\usepackage{microtype}
\geometry{a4paper,margin=3cm}
\setlength{\parindent}{1.25em}
\setlength{\parskip}{0pt}
\emergencystretch=8em
\allowdisplaybreaks
\newcommand{\pair}[2]{\left(#1\middle|#2\right)}
\newcommand{\eps}{\varepsilon}

\begin{document}

\section*{\S\ 5. Тождества разложения в явной форме.}

Если в случае 4) в (25.) специализировать $\tau$ до $\sigma+\rho$, то получается:
\begin{equation}
\begin{aligned}
 \pair{q_\rho A^\sigma}{y^{(1)}\cdots y^{(\rho+\sigma)}}
 &=\sum_{1\le i_1<\cdots<i_\rho\le \rho+\sigma}
 \eps_{i_\rho}\pair{q_\rho}{y^{(i_1)}\cdots y^{(i_\rho)}}\\
 &\qquad\cdot
 \pair{A^\sigma}{y^{(i_{\rho+1})}\cdots y^{(i_{\rho+\sigma})}},
\end{aligned}
\tag{32}
\end{equation}
где $\eps_{i_\rho}=(-1)^{\delta_{i_\rho}}$, $\delta_{i_\rho}=\rho(\rho+1)/2+i_1+\cdots+i_\rho$. Это более общая форма теоремы о произведении для матриц, чем (16.) и (17.); она получается из разложения Лапласа детерминанта продуктовой матрицы
\[
 \sum\pm(v^{(1)}y^{(1)})\cdots(v^{(\rho)}y^{(\rho)})
 (a^{(1)}y^{(\rho+1)})\cdots(a^{(\sigma)}y^{(\rho+\sigma)}).
\]
Следовательно, более общая теорема о произведении составлена из тождеств (20.), соответственно (21.), и этим объясняется выделение специальной формы (16.), (17.).

Соотношение (32.) теперь дает средство для явного представления тождеств разложения. Для этого рассмотрим формы, входящие в (31.), как основные формы; то есть выполним подстановку (11.), которая ничего не меняет в явном представлении через строки $A,B$, поскольку по (8.) вновь введенные строки $R$ выражаются через сами строки $A,B$, а не через их частичные строки $a^{(1)}a^{(2)}\ldots$, $b^{(1)}b^{(2)}\ldots$. Пусть, следовательно,
\begin{equation}
 \pair{q_{\rho-\alpha}B^{n-\rho-\sigma+\tau}}{A_{n-\sigma}p_{\tau-\alpha}}
 =\pair{R_{\rho-\alpha}p_{n-\rho+\alpha}}{R^{\tau-\alpha}p_{\tau-\alpha}},
\tag{33}
\end{equation}
тогда для отдельной суммы из (25.), соответствующей индексу $\alpha$, получаем:
\begin{equation}
\begin{aligned}
&\sum \eps_{i_\alpha}
\pair{R_{\rho-\alpha}p_{n-\rho}y^{(i_1)}\cdots y^{(i_\alpha)}}
{R^{\tau-\alpha}y^{(i_{\alpha+1})}\cdots y^{(i_\tau)}}\\
&\quad=\sum \eps_{i_\alpha}
\pair{[R_{\rho-\alpha}p_{n-\rho}]^\alpha
y^{(i_1)}\cdots y^{(i_\alpha)}}
{R^{\tau-\alpha}y^{(i_{\alpha+1})}\cdots y^{(i_\tau)}}\\
&\quad=\eps\pair{[R_{\rho-\alpha}p_{n-\rho}]^\alpha R^{\tau-\alpha}}{p_\tau}
 =\eps\pair{R^{\tau-\alpha}q_{n-\tau}}{R_{\rho-\alpha}p_{n-\rho}}.
\end{aligned}
\tag{34}
\end{equation}
Тем самым действительно задано явное заключительное выражение, а симметрическая форма, в которой входят $q_\rho$ и $p_\tau$, показывает, что (30.) приводит к тому же заключительному выражению, которое поэтому не зависит от исходного разложения. Различие между (25.) и (30.) существует лишь до тех пор, пока сами строки $y$ и $v$ имеют самостоятельное значение, то есть пока тождество разложения по нашему определению переходит в разложимое тождество. Чтобы перейти к тождествам с большим числом строк символов, нужно по \S\ 3 (конец) лишь разложить строку $q_\rho$ в $q_{\rho_1}C^{\sigma_1}$, соответственно $p_\tau$ в $p_{\tau_1}C_{\tau-\tau_1}$. Тогда возникают выражения:
\begin{equation}
 \pair{q_{\rho_1}C^{\sigma_1}}
 {[R^{\tau-\alpha}q_{n-\tau}]_\lambda R_{\rho_1+\sigma_1-\alpha-\lambda}},
 \qquad
 \pair{[R_{\rho-\alpha}p_{n-\rho}]^\lambda R^{\tau-\alpha}}
 {p_{\tau_1}C_{\tau-\tau_1}},
\tag{35}
\end{equation}
а они имеют точно тот же тип, что и в случае двух строк символов, и потому при подстановке, соответствующей (33.), снова допускают явное представление; повторением получается явное представление для произвольного числа строк. Стоит еще отметить, что первая и последняя суммы в (25.), соответственно (30.), дают явное представление без применения (33.), только по формуле (32.), или вообще сводятся к одному члену, явно содержащему все строки. Соотношение, возникающее из (32.) при разложении $q_\rho$ в $q_{\rho_1}C^{\sigma_1}$ и т. д., можно рассматривать как теорему о произведении в ее самой общей форме.

Подводя итог, покажем:

\emph{Теорема III:} Тождествами
\begin{equation}
 \pair{q_\rho A^\sigma}{p_\tau B_{\rho+\sigma-\tau}}
 =\sum_{\alpha=\max(0,\tau-\sigma)}^{\min(\tau,\rho)}
 \eps\pair{R^{\tau-\alpha}q_{n-\tau}}{R_{\rho-\alpha}p_{n-\rho}},
\tag{36}
\end{equation}
где строки $R$ определяются через (33.), и тождествами, возникающими из них при разложении $q_\rho,p_\tau$ в $q_{\rho_1}C^{\sigma_1}$ и т. д., исчерпывается вся совокупность неразложимых тождеств.

Действительно, возникающие тождества являются самыми общими в своем роде, поскольку отдельные строки уже не подчиняются никаким ограничениям по весу и числу, как это еще было в случае (20.), (21.). Но между этими строками не существует и никаких других тождеств: при разложении строки $p_\tau$ в $y^{(1)}\cdots y^{(\tau)}$ самые общие тождества составлены из (20.), следовательно, по \S\ 3 (конец), из (20a.); причем композиция, обсуждавшаяся после (16.), есть именно та, которая проведена в \S\ 4. Переход к произвольному числу строк, как показывает обсуждение (20a.), дает применение одной и той же операции к выражениям, возникающим при разложении $q_\rho$ в $q_{\rho_1}C^{\sigma_1}$ и т. д. Отсюда также следует, что прямой переход от (20.) вместо (20a.) не дает ничего нового; это легко подтвердить вычислением, если один раз специализировать строку $q_\rho$ в (25.), а другой раз выполнить переход от (20.) методом \S\ 4. В обоих случаях возникают одни и те же суммы, которые при применении общей теоремы о произведении (см. выше) переходят в тождества, вытекающие из (36.). Тождества (20.), (21.) соответствуют специальным случаям $\tau=1$, $\rho=1$; тогда появляются только по две отдельные суммы, а значит, согласно замечанию выше, явное представление без подстановки (33.), что согласуется с прежним.

В заключение кратко упомянем еще два специальных случая тождества разложения. Если в (31.) положить $\tau=\sigma$, $\rho=n-\sigma$ и обозначить строку $p_\tau$ как $p'_\sigma\sim q'_{n-\sigma}$, то получается:
\begin{equation}
\begin{aligned}
\pair{A^\sigma}{p_\sigma}\pair{B^\sigma}{p'_\sigma}
-\pair{A^\sigma}{p'_\sigma}\pair{B^\sigma}{p_\sigma}
=\sum_{\alpha=1}^{\min(\sigma,n-\sigma)}
\Pi\pair{q_{n-\sigma-\alpha}B^\sigma}{A_{n-\sigma}p_{\sigma-\alpha}}.
\end{aligned}
\tag{37}
\end{equation}
Это формула Клебша\footnote{Clebsch, a. a. O., S. 25--32. Клебш доказывает, что отдельные члены правых частей зависят только от строк $A,B$; явное представление получалось бы применением (32.) к его выражениям $\Phi_h$.} для разложения формы с двумя когредиентными строками $p_\sigma,p'_\sigma$ по элементарным ковариантам. Элементарные коварианты, принадлежащие форме $\pair{A^\sigma}{p_\sigma}^m\pair{B^\sigma}{p'_\sigma}^n$, тогда следующие:
\begin{equation}
 \left\{\sum_{\alpha=1}^{\min(\sigma,n-\sigma)}
 \eps\pair{q_{n-\sigma-\alpha}B^\sigma}{A_{n-\sigma}p_{\sigma-\alpha}}
 \right\}^{\rho}
 \pair{A^\sigma}{p_\sigma}^{m-\rho}\pair{B^\sigma}{p'_\sigma}^{n-\rho},
 \qquad (\rho=0,1,\ldots,m,n).
\tag{38}
\end{equation}
Они, как мы коротко будем говорить, возникают из основной формы через ``когредиентную свертку с дефектом $\rho$'' (ср. \S\ 2).

Во-вторых, если положить $\tau=\sigma-\lambda$, $\rho=n-\sigma-\lambda$; $B^\sigma=A^\sigma$, то получается:
\begin{equation}
\begin{aligned}
\pair{q_{n-\sigma-\lambda}A^\sigma}{A_{n-\sigma}p_{\sigma-\lambda}}
&=(-1)^\lambda
\pair{q_{n-\sigma-\lambda}A^\sigma}{A_{n-\sigma}p_{\sigma-\lambda}}\\
&\quad+(-1)^{(n-\sigma)(\sigma-\lambda)}
 \sum_{\alpha=1}^{\min(\sigma-\lambda,n-\sigma-\lambda)}
 \Pi\pair{q_{n-\sigma-\lambda-\alpha}A^\sigma}
 {A_{n-\sigma}p_{\sigma-\lambda}}.
\end{aligned}
\tag{39}
\end{equation}
Следовательно, как отмечено в \S\ 2, условия (12.), характеризующие строки символов, при нечетном дефекте $\lambda$ являются следствием условий более высокого дефекта. Для линейной формы $\pair{A^\sigma}{p_\sigma}=\pair{B^\sigma}{p_\sigma}$ из (39.) следует, что ее неприводимые инвариантные образования, квадратичные по коэффициентам, сводятся к образованиям четного дефекта; это согласуется с известным фактом, что линейная форма в бинарной и тернарной областях не имеет инвариантных образований.

Заметим, что общие тождества первоначально были выведены при предположении, что отдельные строки удовлетворяют условиям (12.). Однако, поскольку эти отдельные строки входят в тождества только линейно, последние справедливы и для более общих строк, не удовлетворяющих условиям (12.).

\end{document}
